\section{Notation and parameters}\label{sec:notation}

All logarithms are natural. Put \(L=\log N\).
The integer \(N\) is even throughout.
Table~\ref{tab:notation} fixes the common notation.

\begingroup
\small
\renewcommand{\arraystretch}{1.15}
\begin{longtable}{@{}p{.27\linewidth}p{.67\linewidth}@{}}
\caption{Common notation}\label{tab:notation}\\
\hline
Symbol & Meaning\\
\hline
\(\mathbb P\) & The set of primes.\\
\(\varphi(n),d(n)\) & Euler's totient and divisor function. The inherited notation \(\tau(n)=d(n)\) has an integer argument; the scalar \(\tau=\sigma^9\) has none.\\
\(\mu_{\mathrm{Mob}}(n),\mu(n)\) & The M\"obius function; use \(\mu_*\) for the sieve cutoff.\\
\(\Omega(n)\) & Prime factors counted with multiplicity.\\
\(P^-(n)\) & Least prime factor; \(P^-(1)=\infty\).\\
\(\mathbf1_E\) & Indicator of \(E\).\\
\(\gamma,\gamma_E\) & Euler's constant, with \(\gamma_E=\gamma\).\\
\(a,b\) & Representation exponent and smaller-factor cutoff.\\
\(\varepsilon\) & Fixed truncation parameter.\\
\(\alpha,\beta,\mu_*,\lambda,t_0,u\) & Fixed sieve exponents in \eqref{eq:not-sieve-parameters}; cutoff \(\mu=\mu_*\).\\
\(C(N),U_C(N),W(N)\) & Normalizations in \eqref{eq:not-normalization}.\\
\(\sigma,\tau,\rho\) & Partition parameters: \(\tau=\sigma^9\), \(\rho=1+\tau\).\\
\(\sigma_*,\tau_*,\rho_*\) & Their fixed choices for the final argument.\\
\(\kappa_{\mathrm{tr}},N_{\mathrm{tr}}\) & Transfer parameters in \eqref{eq:not-transfer-parameters}.\\
\(\delta_{\mathrm{err}}\) & Fixed total error bound, \(10^{-4}\).\\
\(\delta_{\mathrm{pk}}\) & The parameter in the factor-partition theorem.\\
\(h_G,h,\delta_h\) & Profile displacement \(h_G=9/5000\); \(h=h_G\) in \eqref{eq:retB-parameters}, \(\delta_h=h_G/\alpha\).\\
\(\epsilon_s,\delta_s\) & Local sieve parameters, \(\delta_s=\epsilon_s^9\); applications use \(\epsilon_s=\sigma\).\\
\(\eta_L\) & The level margin, \(10^{-4}\); separate from \(\varepsilon\) and \(\delta_{\mathrm{err}}\).\\
\(r,q\) & In \eqref{eq:intro-target}, \(q\) is prime and \(r\) is either \(1\) or prime; other uses are locally bound variables.\\
\(c,I(b)\) & Optional sharper coefficient and the logarithm in \eqref{eq:intro-exterior-constant}.\\
\(c_{K2,G0}\) & Common K2/L3 coefficient in \eqref{eq:intro-fallback-constant}.\\
\(C_E,C_{\mathrm{AA}}\) & Error constants in \eqref{eq:l3-sieve-error} and \eqref{eq:k2-aa-mass}.\\
\(K_{\mathrm{SW}}(A),C_{\mathrm{K2},\sigma},C_{\mathrm{L3},\sigma}\) & The logarithmic Siegel--Walfisz level and the K2 and L3 coefficient bounds; each is local to its stated application.\\
\(R,R_{\mathrm{fix}}\) & The signed sum in \eqref{eq:intro-source}; elsewhere \(R\) is a local level or remainder.\\
\(\Sigma_W,\mathcal W\) & Equal signed source sums; see \eqref{eq:source-weight-identity}.\\
\(W_{13},W_{rs}\) & Arithmetic source weight and spectral operator, respectively. Neither denotes \(W(N)\).\\
\(\overline H,\overline H_L\) & Equal switched-integral bounds, \(389253/5000000\).\\
\(H_{\mathrm{pk,max}},H_{\max}\) & Equal packing ceilings, \(153/256\).\\
\(D_{\mathrm{II}},\overline J_6\) & Original Paper II lower constant and sixth upper endpoint; \eqref{eq:intro-components}, \eqref{eq:ra-j-functional}.\\
\(g_{\mathrm{RA,lo}}\) & Certified floor \(0.0989598353\), used in \eqref{eq:intro-DB}.\\
\(E',E_{\mathrm{van}}\) & Complete retained error and inherited envelope; \eqref{eq:intro-retained-error}, \eqref{eq:retB-base-errors}.\\
\(\begin{gathered}\Lambda_{j,N},\\\Lambda_{G,N}\end{gathered}\) & Retained conversion differences, defined in \eqref{eq:retB-direct-Lambda} and \eqref{eq:retB-G-array}.\\
\(\Delta_G,\Lambda_{G,N}\) & Joint-profile gain and conversion difference in \eqref{eq:retB-G-array}. The spelling \(G1\) in subscripts denotes the same group.\\
\(\begin{gathered}E_{\mathrm{RA}},E_{\mathrm{sw}},\\E_K,E_{\mathrm{exc}}\end{gathered}\) & Nonnegative complete errors in \eqref{eq:intro-comparisons}--\eqref{eq:intro-error-identifications} and \eqref{eq:intro-exterior}.\\
\(\mathcal E_B\) & Their total, including \(E'\), in \eqref{eq:intro-total-error}.\\
\(D_B\) & The constant in \eqref{eq:intro-DB}; evaluated in \eqref{eq:bill-constant}.\\
\(\mathcal A_N,\mathcal P_b,G_b,B_b\) & Shifted-prime source, admissible factors, positive count and exterior count in \eqref{eq:intro-count}, \eqref{eq:intro-source}.\\
\(S_1,\ldots,S_{13},S_{\mathrm{sw}}\) & Paper II arithmetic sums retained in \eqref{eq:intro-source}; definitions in Section~\ref{sec:source-decomposition}.\\
\(\mathfrak g_{\mathrm{RA}},g_{\mathrm{RA,lo}}\) & Integral gain and its smaller rational floor; only the floor enters $D_B$.\\
\(\begin{gathered}g_6,M_{5,\mathrm{internal}},\\K_{\mathrm{hi}},L_{G1}\end{gathered}\) & Optional six-endpoint gain, shifted-profile gain and mass constants; the main proof sets \(g_6=0\), see \eqref{eq:retB-six-gain}, \eqref{eq:intro-components}, \eqref{eq:intro-G1-constants}.\\
\(\Delta D_{L3}\) & Switched gain $(2-c_{K2,G0}/4)\overline H$.\\
\(G_*,G_0,U_0,c_{\mathrm{O1}}\) & Common unshifted K2/L3 functional; coefficient cap $G_0=6088440363/10^9$; optional sharper upper certificate $U_0$; its exact coefficient in \eqref{eq:l3-exact-certificate-constant}.\\
\(\mathcal H,\Lambda^0_{j,N}\) & Retained row set and saved-array contribution in \eqref{eq:retB-goal}.\\
\(\mathcal E,E_{\mathrm{pt}}\) & Pointwise exceptional set and error in \eqref{eq:source-exceptional-set}, \eqref{eq:source-pointwise-error}.\\
\(N_0^*,N_{\mathrm{floor}},T_{\bullet}\) & Common threshold, numerical floor and named threshold requirements in \eqref{eq:not-common-threshold}.\\
\(\zeta,a_{\mathrm{buf}}^*,\mathfrak E_*\) & Buffer $10^{-6}$, scalar endpoint and its exponential in \eqref{eq:bill-critical-exponent}.\\
\(W_{\mathrm{AA},K},W_*\) & K2 source-weight bound in \eqref{eq:k2-aa-mass} and Paper II bound $10^{14}$; these are distinct.\\
\(A_K,P_K,r_K\) & K2 error multiplier, peel multiplier and vanishing remainder in \eqref{eq:bill-K-error}.\\
\(C_{\mathrm{sw}},P_{\mathrm{sw}},V_{\mathrm{sw}}\) & Switched error multipliers and positive remainder in \eqref{eq:l3-price}, \eqref{eq:l3-remainder}.\\
\(L,g_{\mathrm{AA}}\) & $\log N$ and Paper II's fixed error allowance in \eqref{eq:bill-paperII-AA}.\\
\hline
\end{longtable}
\endgroup

Cusps carry the index \(\mathfrak a\) in spectral sums.
The switched blocks carry the index \(\mathfrak b\).
Partition cells are \(c\) or \(\mathfrak c\).
Coefficient sequences are \(a_n,b_m\), with dummy indices interchanged when needed.
Spectral moduli and cusp numerators are locally bound variables.
They do not redefine the constants \(c\) and \(u\).
Packing coordinates also have local scope.
Bare \(T\) and \(w\) are local endpoints or weights; the named saved and
root-A scales are \(T_{\mathrm{save}},T_{\mathrm{int}}\), and \(w_{\mathrm{RA}}\).
The exponent \(A_0>0\) in distribution estimates is a logarithmic precision.
The normalization factor is written \(2e^{-\gamma}\).

\subsection{Fixed parameters}

Keep \(a,b,\varepsilon\) from \eqref{eq:intro-exponents}.
The sieve parameters are
\begin{equation}\label{eq:not-sieve-parameters}
\begin{aligned}
 \alpha&=\frac{100}{1149},&\beta&=\frac{50}{309},&
 \mu_*&=\frac{25}{128},\\
 \lambda&=\frac{93}{400},&t_0&=\frac{57}{224},&
 u&=\frac1{500}.
\end{aligned}
\end{equation}
These follow Paper II's notation
\cite[(1.5)]{PaperII}; $\varepsilon$ is $\varepsilon_A$ in
\cite[(1.3)]{PaperII}.
Put
\begin{equation}\label{eq:not-normalization}
\begin{aligned}
 C_2&=\prod_{\substack{\ell\in\mathbb P\\\ell>2}}
       \left(1-\frac1{(\ell-1)^2}\right),\\
 C(N)&=C_2\prod_{\substack{\ell\in\mathbb P,\ \ell\mid N\\\ell>2}}
                   \frac{\ell-1}{\ell-2},\\
 U_C(N)&=\frac{C(N)N}{L^2},\qquad
 W(N)=4(1-\varepsilon)U_C(N).
\end{aligned}
\end{equation}
We often suppress the argument \(N\).

Define the transfer parameters by
\begin{equation}\label{eq:not-transfer-parameters}
\begin{aligned}
 \kappa_{\mathrm{tr}}&=a(1-b)-1=\frac7{4000000},\\
 N_{\mathrm{tr}}&=\varepsilon^{-(a-1)/\kappa_{\mathrm{tr}}}
                 =\varepsilon^{-3000000/7}.
\end{aligned}
\end{equation}
Here $a-1=3/4$ is the exponent on the larger prime in the bound for the
smaller factor.
The transfer argument appears in the final section.

\subsection{Order of choice}

First fix all \(\sigma\)-independent constants in the estimates.
These include \(C_E\) in \eqref{eq:l3-sieve-error}
and \(C_{\mathrm{AA}}\) in \eqref{eq:k2-aa-mass}.
Then choose one rational \(\sigma_*\), with
\begin{equation}\label{eq:not-sigma-choice}
 0<\sigma_*\le10^{-10},\qquad
 \tau_*=(\sigma_*)^9,\qquad \rho_*=1+\tau_*.
\end{equation}
The final section gives the remaining upper restrictions.
Apply every partition estimate with \(\sigma=\sigma_*\).
Freeze its finite partitions and coefficient bounds next.
Freeze the resulting finite-family bounds and logarithmic precisions too.
These choices precede the common threshold.

Use the following named thresholds. Each requires its listed statement
for every even $N$ above it, at the fixed $\sigma_*$. The cited comparison
propositions establish these statements. The numerical floor supplies none
of them. The resulting threshold exists but may be ineffective.
\begin{enumerate}
\item $T_{\mathrm{src}}$: the actual-source identity and first line of
\eqref{eq:intro-comparisons}, including the pointwise hypotheses in
\eqref{eq:source-pointwise} and \eqref{eq:source-first-lower}.
Include the explicit pointwise threshold
$N_{\mathrm{pt}}=1+\lceil\max\{2^{1/\alpha},\varepsilon^{-3},
\varepsilon^{-1/(1-\lambda-\beta)}\}\rceil$ from
Section~\ref{sec:source-decomposition}.
\item $T_{\mathrm{ret}}$: the retained comparison, saved-array
identifications and nonnegativity of $E_{\mathrm{van}}$ in
Section~\ref{sec:retained-side}.
\item $T_{\mathrm{RA}}$: the sixth-sum comparison
\eqref{eq:ra-b11-floor}, with every distribution and coefficient bound
specified in Section~\ref{sec:root-a}. It also covers the finite source
threshold $N_6$, chosen so that every even $N\ge N_6$ satisfies
\[
 N^{1/6}>L^{\max(B_L,B_{\mathrm{tail}})},\quad
 N^{v_{\mathrm{RA}}-1/3}>L^B,\quad
 \varepsilon N^{1/2}>N^{\mu_*},\quad
 NL^{-333}\ge N^{\mu_*}.
\]
Thus the operative $T_{\mathrm{RA}}$ is the maximum of the sixth-sum
comparison threshold and this $N_6$.
\item $T_K$: the actual-source exterior comparison
\eqref{eq:k2-final-fixed-sigma} and its nonnegative error decomposition,
the $r_K$ formula displayed after
\eqref{eq:k2-final-fixed-sigma}.
\item $T_{\mathrm{sw}}$: the switched comparison \eqref{eq:l3-final}
and the nonnegative decomposition \eqref{eq:l3-remainder}.
\item $T_{\mathrm{len}}$: all scale-conversion, zero-mode, factorization,
packing, interval-length and logarithmic-precision conditions of the
preceding applications, with the fixed finite partitions. In particular,
require $N^{\sigma_*}\ge2$,
$N^{60\sigma_*2^{-126}}\ge12L^{334}$,
\eqref{eq:k2-med-length} and \eqref{eq:l3-lengths}, and include
$N_{\mathrm{dom}},N_{z,0},N_{\mathrm{LEN,app}}$ from
\eqref{eq:l3-structural-thresholds}. The remaining application conditions
are exactly those of the cited comparison propositions. The floor implies
none of them.
\item $T_{E'},T_V,T_{E_{\mathrm{RA}}},T_r,T_{E_{\mathrm{exc}}}$:
respectively the eventual bounds
\[
\begin{gathered}
 E'\le\frac{\delta_{\mathrm{err}}}{20},\qquad
 V_{\mathrm{sw}}\le\frac{\delta_{\mathrm{err}}}{20},\qquad
 E_{\mathrm{RA}}\le\frac{\delta_{\mathrm{err}}}{10},\\
 r_K\le\frac{\delta_{\mathrm{err}}}{10},\qquad
 E_{\mathrm{exc}}\le\frac{\delta_{\mathrm{err}}}{10}.
\end{gathered}
\]
All arguments are $(N;\sigma_*)$, or $(N;b,\sigma_*)$ for $r_K$.
\end{enumerate}
These are finitely many requirements after all finite choices.
Define the common threshold once by
\begin{equation}\label{eq:not-common-threshold}
\begin{aligned}
 \mathcal T_*={}&\{T_{\mathrm{src}},T_{\mathrm{ret}},T_{\mathrm{RA}},
 T_K,T_{\mathrm{sw}},T_{\mathrm{len}},\\
 &\hspace{12mm}T_{E'},T_V,T_{E_{\mathrm{RA}}},T_r,T_{E_{\mathrm{exc}}}\},\\
 N_0^*(\sigma_*)&=2\left\lceil\frac12
 \max\bigl(\{N_{\mathrm{tr}},N_{\mathrm{floor}}\}\cup\mathcal T_*\bigr)
 \right\rceil,\\
 N_{\mathrm{floor}}&=2\left\lceil e^{2^{150}}/2\right\rceil.
\end{aligned}
\end{equation}
The floor is a convention added to Paper~I's
selected-interval length hypothesis
\cite[\S3.2, length hypothesis preceding (3.12)]{PaperI}. It proves none of
the listed estimates. The named application estimates establish this
threshold. The resulting $N_0^*$ may be ineffective.

This implements the order in \eqref{eq:intro-order}.
Afterwards, \(N\) is even and $N\ge N_0^*(\sigma_*)$.

Let \(\epsilon_{\mathrm{err}}(N;\sigma_*)\) be a real- or complex-valued error function.
Write \(\epsilon_{\mathrm{err}}(N;\sigma_*)=o_{\sigma_*}(1)\) if and only if
\begin{equation}\label{eq:not-little-o}
\begin{gathered}
 \forall\eta>0\ \exists N_\eta(\sigma_*)\ \forall N\in2\mathbb N,\\
 N\ge N_\eta(\sigma_*)\ \Longrightarrow\ |\epsilon_{\mathrm{err}}(N;\sigma_*)|\le\eta.
\end{gathered}
\end{equation}
All preceding fixed data are suppressed here.
The parameter \(\sigma_*\) remains fixed in this limit.
Uniformity in other variables is stated with each estimate.
